\chapter{Alexander Grothendieck}\label{ch1966alexandergrothendieck:alexander-grothendieck-fields-medal-1966-a-comprehensive-report}

\begin{center}\rule{0.5\linewidth}{0.5pt}\end{center}

\textbf{Prepared:} July 2026
\textbf{Subject:} Alexander Grothendieck (b. March 28, 1928, Berlin, Germany; d. November 2014, Saint-Girons, France)
\textbf{Position:} Professor, IHÉS; University of Montpellier
\textbf{Doctoral Advisor:} Laurent Schwartz (Université Henri Poincaré, Nancy)
\textbf{Fields Medal Citation:} \emph{``Built on work of Weil and Zariski and
effected fundamental advances in algebraic geometry. He introduced the idea of
K-theory. Revolutionized homological algebra in his well-known `Tohoku paper'.''}

\begin{center}\rule{0.5\linewidth}{0.5pt}\end{center}

\begin{prereqs} Abstract algebra (rings, ideals, modules), basic category theory
(functors, natural transformations), point-set topology.
\end{prereqs}

\begin{objectives} Explain the scheme-theoretic generalization of algebraic
geometry \cite{grothendieck1960}; describe étale cohomology; understand motives as universal cohomology
theories; introduce Grothendieck topologies and descent theory.
\end{objectives}

\section{Problem History}

In the early 20th century, algebraic geometry rested on classical varieties over
algebraically closed fields and Weil's \emph{Foundations} (1946). The theory was
built around the maximal spectrum of a ring and Zariski's notion of a variety as
an irreducible closed subset of projective space. While powerful, this framework
had several serious limitations.

A fundamental problem was the \textbf{Weil conjectures} (1949): for a smooth
projective variety $X$ over $\mathbb{F}_q$, the zeta function
\[
Z(X, t) = \exp\left(\sum_{n=1}^{\infty} |X(\mathbb{F}_{q^n})| \frac{t^n}{n}\right)
\]
should be rational, satisfy a functional equation, and obey a Riemann hypothesis.
These conjectures required a cohomology theory in characteristic $p$ that
behave like singular cohomology. The Weil conjectures provided a unifying vision
that connected number theory, topology, and algebraic geometry in an
unexpected way.

The Zariski topology is too coarse for useful cohomology. In this topology, open
sets are complements of algebraic subvarieties, and the resulting sheaf
cohomology of coherent sheaves vanishes for $i > 0$ on affine schemes by
Serre's criterion. This vanishing means the Zariski topology cannot capture the
global topological information needed for the Weil conjectures. Weil proposed a
``topology'' using correspondences, but this lacked the axiomatic framework needed
for computations and generalizations.

\subsection*{The Failure of Classical Topologies}

The étale fundamental group of $\operatorname{Spec}(\mathbb{F}_q)$ is $\hat{\mathbb{Z}}$,
the profinite completion of $\mathbb{Z}$, which captures the Galois group of the
algebraic closure. This suggested that a suitable ``topology'' on algebraic
varieties should mimic the role of covering spaces in classical topology, but
with finite étale morphisms replacing covering maps. The Zariski topology's
coarse nature---where every non-empty open set is dense---meant that sheaf
cohomology in this topology could not detect the subtleties needed for a
Weil cohomology theory.

Homological algebra (Cartan--Eilenberg, 1956) provided derived functors and
spectral sequences but lacked a unified categorical setting. Sheaf theory worked
on topological spaces but needed algebraic generalization. Grothendieck's
approach was to rebuild algebraic geometry from categorical foundations,
replacing varieties with \textbf{schemes}---spaces locally modeled on prime
spectra of commutative rings.

The problem of finding a Weil cohomology theory---satisfying Poincaré duality,
the Lefschetz trace formula, and the Künneth formula---drove Grothendieck's
work on étale cohomology. Without such a theory, the Weil conjectures could not
even be formulated precisely. The key insight was that one needed to replace the
Zariski topology with a finer topology that retained enough algebraic structure
while allowing meaningful cohomological computations.

\subsection*{The Search for a Weil Cohomology}

Serre had already developed the theory of coherent sheaves and their cohomology
on projective varieties, establishing that for a coherent sheaf $\mathcal{F}$ on
a projective variety $X$, the cohomology groups $H^i(X, \mathcal{F})$ are
finite-dimensional. However, coherent cohomology vanishes on affine varieties,
making it useless for the topological purposes required. The quest was for a
cohomology theory that would satisfy the following properties:

\begin{enumerate}
\item Finiteness: $H^i(X, \mathbb{Q}_\ell)$ should be finite-dimensional.
\item Poincaré duality: There should be a perfect pairing
$H^i(X) \times H^{2n-i}(X) \to \mathbb{Q}_\ell(-n)$.
\item Lefschetz trace formula: The number of fixed points of a morphism
should be expressible as an alternating sum of traces on cohomology.
\item Künneth formula: $H^*(X \times Y) \cong H^*(X) \otimes H^*(Y)$.
\end{enumerate}

\section{How Was It Solved}

Grothendieck's revolution began with the \textbf{Tohoku paper} (1957),
establishing modern homological algebra. He introduced \textbf{abelian
categories} with axioms AB1--AB5, proving that any AB5 category with a generator
has enough injectives. This foundational result allowed the systematic use of
derived functors in contexts far beyond modules over rings.

\subsection*{Abelian Categories and Derived Functors}

An \textbf{abelian category} $\mathcal{A}$ is a category with a zero object,
biproducts, kernels, cokernels, and where every monomorphism is a kernel and
every epimorphism is a cokernel. Grothendieck introduced the axioms:

\begin{itemize}
\item \textbf{AB1}: Every family of subobjects has an intersection.
\item \textbf{AB2}: The canonical factorization $f = me$ of any morphism has
$e$ epimorphic and $m$ monomorphic.
\item \textbf{AB3}: Every family of objects has a direct sum.
\item \textbf{AB4}: Direct sums preserve epimorphisms.
\item \textbf{AB5}: Filtered colimits are exact.
\item \textbf{AB6}: Filtered colimits commute with finite limits.
\end{itemize}

The \textbf{Tohoku theorem} states that in an AB5 category with a generator,
every object admits an injective resolution, and the right derived functors of
any left exact functor exist. This established the formal framework for
cohomology theories in algebra and geometry.

\subsection*{Scheme Theory}

A \textbf{scheme} is a locally ringed space $(X, \mathcal{O}_X)$ where every
point has a neighborhood isomorphic to $\operatorname{Spec}(R)$ for some
commutative ring $R$. Here $\operatorname{Spec}(R)$ denotes the set of prime
ideals of $R$ equipped with the Zariski topology and the structure sheaf.

The category of schemes $\mathbf{Sch}$ admits fibered products $X \times_S Y$
and morphisms (flat, proper, étale). The \textbf{scheme} $\operatorname{Spec}(\mathbb{Z})$
contains both characteristic $p$ and characteristic $0$ information, providing
a unified framework for arithmetic geometry.

Grothendieck introduced the \textbf{functor of points} viewpoint: a scheme $X$
over $S$ is determined by the functor
\[
X(S') = \operatorname{Hom}_S(S', X)
\]
for all $S$-schemes $S'$. This Yoneda-style approach allows one to study
schemes through their points over all base changes, which is essential for
moduli theory and for understanding deformations.

The \textbf{EGA} (Éléments de géométrie algébrique), co-authored with Jean
Dieudonné, provided the foundational texts on scheme theory. Key constructions
include:
\begin{itemize}
\item The \textbf{structure morphism} $f: X \to S$ for a scheme over a base.
\item The \textbf{relative Spec} construction $\operatorname{Spec}(\mathcal{A})$
for a quasicoherent sheaf of algebras $\mathcal{A}$.
\item The \textbf{relative Proj} construction $\operatorname{Proj}(\mathcal{S})$
for a graded algebra sheaf.
\item \textbf{Grothendieck's existence theorem} relating formal schemes to
algebraic schemes.
\end{itemize}

\subsection*{Étale Cohomology}

The \textbf{étale topology} uses étale morphisms as coverings. A morphism
$f: X \to Y$ is \textbf{étale} at $x \in X$ if it is flat and unramified,
meaning the sheaf of relative differentials $\Omega_{X/Y}$ vanishes at $x$.
Étale morphisms are the algebraic analogues of local homeomorphisms in
topology.

For a scheme $X$ and a torsion sheaf $\mathcal{F}$ on the étale site $X_{\text{ét}}$,
the \textbf{étale cohomology groups} $H^i_{\text{ét}}(X, \mathcal{F})$ are
computed as the right derived functors of the global section functor. The
comparison theorem states that for a smooth variety $X$ over $\mathbb{C}$,
there is a canonical isomorphism
\[
H^i_{\text{ét}}(X, \mathbb{Z}/n\mathbb{Z}) \cong H^i(X^{\text{an}}, \mathbb{Z}/n\mathbb{Z})
\]
for $n$ invertible on $X$, connecting étale cohomology to classical singular
cohomology.

For a prime $\ell$ invertible on $X$, one defines the \textbf{$\ell$-adic
cohomology} groups:
\[
H^i_{\text{ét}}(X, \mathbb{Q}_\ell) = H^i_{\text{ét}}(X, \mathbb{Z}_\ell) \otimes_{\mathbb{Z}_\ell} \mathbb{Q}_\ell
\]
where $\mathbb{Z}_\ell = \varprojlim_n \mathbb{Z}/\ell^n\mathbb{Z}$.

The \textbf{Weil conjectures} were proved step by step:
\begin{enumerate}
\item \textbf{Rationality} (Dwork, 1960): $Z(X,t)$ is a rational function.
\item \textbf{Functional equation} (Grothendieck--Deligne): Follows from
Poincaré duality in étale cohomology.
\item \textbf{Riemann hypothesis} (Deligne, 1974): The eigenvalues of
Frobenius on $H^i$ have absolute value $q^{i/2}$.
\end{enumerate}

\subsection*{The Lefschetz Trace Formula}

For a smooth projective variety $X$ over $\mathbb{F}_q$ and the Frobenius
morphism $F: X \to X$, the \textbf{Lefschetz trace formula} states:
\[
|X(\mathbb{F}_q)| = \sum_{i=0}^{2n} (-1)^i \operatorname{tr}\left(F^* | H^i_{\text{ét}}(X_{\overline{\mathbb{F}}_q}, \mathbb{Q}_\ell)\right)
\]
This formula relates the arithmetic count of rational points to the
topological data of étale cohomology. The functional equation of the zeta
function follows from Poincaré duality, which provides a perfect pairing:
\[
H^i_{\text{ét}}(X, \mathbb{Q}_\ell) \times H^{2n-i}_{\text{ét}}(X, \mathbb{Q}_\ell(n)) \to \mathbb{Q}_\ell
\]

\subsection*{Grothendieck--Riemann--Roch and K-theory}

The \textbf{Grothendieck--Riemann--Roch theorem} relates the Chern character of
the direct image of a coherent sheaf to the direct image of the product of the
Chern character and the Todd class:
\[
\operatorname{ch}(f_!(E)) \cdot \operatorname{td}(T_Y) = f_*(\operatorname{ch}(E) \cdot \operatorname{td}(T_X))
\]
where $f: X \to Y$ is a proper morphism of smooth varieties, $E$ is a vector
bundle on $X$, and $f_!$ denotes the derived pushforward in K-theory.

This theorem introduced \textbf{K-theory} $K(X)$---the Grothendieck group of
vector bundles on $X$, defined as the quotient of the free abelian group on
isomorphism classes of vector bundles by the relation $[E] = [E'] + [E'']$ for
every short exact sequence $0 \to E' \to E \to E'' \to 0$.

\subsection*{Descent Theory}

\textbf{Descent theory} addresses the question: when does a sheaf or vector
bundle on a covering $Y \to X$ come from a sheaf on $X$? Grothendieck
formulated this in terms of \textbf{descend data}: for a faithfully flat
morphism $f: Y \to X$, a sheaf $\mathcal{F}_Y$ on $Y$ descends to $X$ if and
only if there is an isomorphism $\phi: p_1^*\mathcal{F}_Y \xrightarrow{\sim}
p_2^*\mathcal{F}_Y$ on $Y \times_X Y$ satisfying the cocycle condition on
$Y \times_X Y \times_X Y$.

The fundamental theorem of descent states that for a faithfully flat
morphism $f: Y \to X$, the category of quasi-coherent $\mathcal{O}_X$-modules
is equivalent to the category of quasi-coherent $\mathcal{O}_Y$-modules with
descent data. This is a cornerstone of modern algebraic geometry and has
profound applications to moduli problems.

\section{Impact of This Discovery}

Grothendieck's schemes became the language of modern algebraic geometry, used in
the proof of Fermat's Last Theorem and the Weil conjectures.

\subsection*{Arithmetic Geometry}

Étale cohomology opened arithmetic geometry and made the Weil conjectures
tractable. The proof of Fermat's Last Theorem by Andrew Wiles (1994) relied
crucially on the language of schemes and étale cohomology developed by
Grothendieck. The modularity theorem, which implies Fermat's Last Theorem, is
stated in terms of elliptic curves over $\mathbb{Q}$ viewed as schemes.

The \textbf{Brauer group} $\operatorname{Br}(X) = H^2_{\text{ét}}(X, \mathbb{G}_m)$
of a scheme plays a central role in arithmetic, controlling obstructions to
the Hasse principle. The \textbf{Birch--Swinnerton-Dyer conjecture}, one of the
Millennium Prize Problems, relates the rank of an elliptic curve (viewed as a
scheme) to the order of vanishing of its L-function at $s=1$.

\subsection*{Motives and Universal Cohomology}

The theory of \textbf{motives} guides modern algebraic geometry. Grothendieck
conceived motives as universal objects from which all cohomology theories arise
by realization functors. The category of motives uses correspondences as
morphisms.

A \textbf{motive} over a field $k$ should be an object $M(X)$ associated to a
variety $X$ such that for any Weil cohomology theory $h^*$, there is a
realization functor $r_h: \mathbf{Mot} \to \mathbf{Vect}$ with
$r_h(M(X)) \cong h^*(X)$. The category of motives is expected to be abelian
with a Tannakian structure, though this remains conjectural.

Voevodsky's triangulated categories of motives proved the \textbf{Milnor
conjecture} (Fields Medal 2002), establishing an isomorphism between
Milnor K-theory and Galois cohomology:
\[
K_n^M(k)/2 \cong H^n_{\text{ét}}(k, \mathbb{Z}/2\mathbb{Z})
\]

\subsection*{Homological Algebra}

The Tohoku paper's framework of abelian categories and derived functors is used in
algebraic topology, algebraic geometry, representation theory, and mathematical
physics. The language of derived categories, introduced by Grothendieck and
developed by Verdier, has become indispensable in modern mathematics.

\textbf{Topos theory} unified topology, algebraic geometry, and logic, with
applications to type theory and computer science. A \textbf{topos} is a category
that behaves like the category of sheaves on a topological space. The
\textbf{classifying topos} for a geometric theory provides a semantic framework
for logic.

The \textbf{Grothendieck--Riemann--Roch theorem} is a workhorse of algebraic
geometry, leading to Quillen's algebraic K-theory (Fields Medal 1978). Quillen
defined higher K-groups $K_n(R)$ using the $Q$-construction, generalizing the
classical $K_0$ to all dimensions.

\subsection*{Dessins d'Enfants}

\textbf{Dessins d'enfants} (children's drawings) provide a concrete description
of the Galois action on combinatorial objects. Grothendieck introduced these
bipartite graphs embedded on Riemann surfaces to study the absolute Galois
group $\operatorname{Gal}(\overline{\mathbb{Q}}/\mathbb{Q})$. Each dessin
corresponds to a Belyi map $f: C \to \mathbb{P}^1$ ramified only over $\{0,
1, \infty\}$, and the Galois group acts faithfully on the set of dessins.

\section{Modern Views and Improvements}

\subsection*{Motivic Homotopy Theory}

Motives have seen major progress through \textbf{motivic homotopy theory},
developed by Morel and Voevodsky. The category $\mathbf{Spc}(k)$ of motivic
spaces is a homotopy-theoretic framework that recovers classical results while
providing new computational tools. The \textbf{motivic stable homotopy category}
$\mathbf{SH}(k)$ provides a unified setting for various cohomology theories.

The \textbf{A1-homotopy theory} replaces the interval $[0,1]$ with the
affine line $\mathbb{A}^1$, leading to a homotopy theory over arbitrary
schemes. This theory has deep connections to algebraic K-theory and
motivic cohomology.

\subsection*{Étale Cohomology and the Six Operations}

Étale cohomology extends to $\ell$-adic sheaves and the \textbf{six operations}
formalism (Kashiwara--Schapira):
\begin{center}
$f^*$, $f_*$, $f_!$, $f^!$, $\otimes$, $\mathcal{H}om$
\end{center}
These operations satisfy adjunctions, base change, and projection formulas that
generalize classical topological results. \textbf{Perverse sheaves} are central
to the geometric Langlands program.

\subsection*{The Geometric Langlands Program}

The Langlands program is deeply rooted in Grothendieck's ideas. The
\textbf{geometric Langlands conjecture} relates D-modules on the moduli stack
of bundles to sheaves on the moduli stack of local systems. This program has
found surprising connections to physics through mirror symmetry and
topological field theory.

\subsection*{Perfectoid Spaces}

\textbf{Perfectoid spaces}, introduced by Peter Scholze, extend scheme theory to
characteristic $p$ in a way that allows comparison with characteristic $0$.
Scholze's proof of the \textbf{weight monodromy conjecture} and the
\textbf{local Langlands correspondence for $\mathrm{GL}_n$} used perfectoid spaces,
building on Grothendieck's foundational work. Perfectoid spaces have become a
central tool in arithmetic geometry.

\subsection*{Condensed Mathematics}

\textbf{Condensed mathematics} (Clausen--Scholze) provides a new foundation for
analytic and algebraic geometry. The category of condensed abelian groups
replaces topological abelian groups and resolves technical issues with
completions and products. This framework has applications to the study of
Banach spaces and p-adic geometry.

\subsection*{The Standard Conjectures}

The \textbf{Standard Conjectures} on algebraic cycles remain open. These
conjectures predict the existence of a Weil cohomology theory with good
properties:
\begin{itemize}
\item \textbf{Conjecture A}: Homological equivalence equals algebraic equivalence.
\item \textbf{Conjecture B}: The Lefschetz operator is invertible on cohomology.
\item \textbf{Conjecture C}: There exists a ``hard Lefschetz'' isomorphism.
\end{itemize}
These conjectures would imply the Hodge conjecture and provide a deeper
understanding of algebraic cycles.

\subsection*{Modern Developments}

Active areas include motives for non-algebraic varieties, the Standard
Conjectures on algebraic cycles, and condensed mathematics (Clausen--Scholze)
as a framework for algebraic geometry in positive characteristic. The
\textbf{motivic integration} program, developed by Kontsevich and others,
provides powerful tools for studying singularities and mirror symmetry.

\begin{center}\rule{0.5\linewidth}{0.5pt}\end{center}

\section*{Summary}
\addcontentsline{toc}{section}{Summary}

Grothendieck revolutionized algebraic geometry by introducing schemes, étale
cohomology, and the language of categories. His vision of motives as universal
cohomology theories and his foundations in the EGA and SGA séminaires reshaped the
field and continue to guide modern research in arithmetic geometry and number
theory. The development of derived categories, descent theory, and K-theory
provided new tools that permeate modern mathematics. Perfectoid spaces and
condensed mathematics represent the latest developments in Grothendieck's
program of understanding geometry through categorical and cohomological lenses.
